\section{Singly Linked List: Construction and Basic Operations}
\label{sec:ll-singly-basics}

\problemstatement{Implement the foundational operations on a singly linked
list: build from an array, traverse, find the length, search for a value,
insert at the head or tail, and delete a value.}

\begin{patternbox}
No pattern recognition is needed here -- this is the machinery every later
section rests on. The one habit to internalise is keeping a \texttt{prev}
handle while walking, so a node can be unlinked or a new node spliced in
without losing the rest of the chain.
\end{patternbox}

\subsection*{Walking with a pointer}

Each node stores a value and a pointer to the next node; the last node
points to \texttt{nullptr}. Traversal is a loop that follows \texttt{next}
until null. Insertion at the head is $O(1)$ (point the new node at the old
head). Insertion at the tail and deletion require walking to the right spot,
keeping the previous node so its \texttt{next} can be rewired.

\begin{ideabox}[Keep the Previous Node When You Might Rewire]
Deletion and mid-list insertion both need to change the \texttt{next} of the
node \emph{before} the target. Tracking \texttt{prev} as you traverse makes
these $O(1)$ once the position is found, with no re-walk.
\end{ideabox}

\subsection*{C++ implementation}

\begin{lstlisting}
#include <bits/stdc++.h>
using namespace std;

struct ListNode {
    int val;
    ListNode* next;
    ListNode(int v) : val(v), next(nullptr) {}
};

ListNode* build(const vector<int>& a) {           // build from array
    ListNode dummy(0), *tail = &dummy;
    for (int x : a) { tail->next = new ListNode(x); tail = tail->next; }
    return dummy.next;
}

int length(ListNode* head) {
    int n = 0;
    for (ListNode* c = head; c; c = c->next) n++;
    return n;
}

bool search(ListNode* head, int key) {
    for (ListNode* c = head; c; c = c->next) if (c->val == key) return true;
    return false;
}

ListNode* insertHead(ListNode* head, int v) {
    ListNode* node = new ListNode(v);
    node->next = head;
    return node;                                  // new head
}

ListNode* deleteValue(ListNode* head, int key) {
    ListNode dummy(0); dummy.next = head;
    ListNode* prev = &dummy;
    while (prev->next && prev->next->val != key) prev = prev->next;
    if (prev->next) prev->next = prev->next->next;   // unlink
    return dummy.next;
}

void print(ListNode* head) {
    for (ListNode* c = head; c; c = c->next) cout << c->val << " ";
    cout << "\n";
}

int main() {
    ListNode* head = build({1, 2, 3, 4});
    print(head);                                  // 1 2 3 4
    cout << "length " << length(head) << ", search 3? " << search(head, 3) << "\n";
    head = insertHead(head, 0);
    head = deleteValue(head, 3);
    print(head);                                  // 0 1 2 4
    return 0;
}
\end{lstlisting}

\begin{complexitybox}
\textbf{Time Complexity.} $O(1)$ head insert; $O(n)$ for length, search,
tail insert, and delete. \textbf{Space Complexity.} $O(1)$ extra (plus the
list itself).
\end{complexitybox}

\begin{takeawaybox}
The dummy node plus a \texttt{prev} handle removes the ``deleting the head''
special case and is the scaffolding for nearly every operation in this
chapter. Master this template before the pointer-heavy problems ahead.
\end{takeawaybox}
